% Options for packages loaded elsewhere
% Options for packages loaded elsewhere
\PassOptionsToPackage{unicode}{hyperref}
\PassOptionsToPackage{hyphens}{url}
\PassOptionsToPackage{dvipsnames,svgnames,x11names}{xcolor}
%
\documentclass[
  spanish,
  10pt,
  letterpaper,
  twocolumn]{article}
\usepackage{xcolor}
\usepackage[top=1.0cm,bottom=1.5cm,left=1.5cm,right=1.5cm]{geometry}
\usepackage{amsmath,amssymb}
\setcounter{secnumdepth}{-\maxdimen} % remove section numbering
\usepackage{iftex}
\ifPDFTeX
  \usepackage[T1]{fontenc}
  \usepackage[utf8]{inputenc}
  \usepackage{textcomp} % provide euro and other symbols
\else % if luatex or xetex
  \usepackage{unicode-math} % this also loads fontspec
  \defaultfontfeatures{Scale=MatchLowercase}
  \defaultfontfeatures[\rmfamily]{Ligatures=TeX,Scale=1}
\fi
\usepackage{lmodern}
\ifPDFTeX\else
  % xetex/luatex font selection
  \setmainfont[]{Fira Sans}
\fi
% Use upquote if available, for straight quotes in verbatim environments
\IfFileExists{upquote.sty}{\usepackage{upquote}}{}
\IfFileExists{microtype.sty}{% use microtype if available
  \usepackage[]{microtype}
  \UseMicrotypeSet[protrusion]{basicmath} % disable protrusion for tt fonts
}{}
\usepackage{setspace}
\makeatletter
\@ifundefined{KOMAClassName}{% if non-KOMA class
  \IfFileExists{parskip.sty}{%
    \usepackage{parskip}
  }{% else
    \setlength{\parindent}{0pt}
    \setlength{\parskip}{6pt plus 2pt minus 1pt}}
}{% if KOMA class
  \KOMAoptions{parskip=half}}
\makeatother
% Make \paragraph and \subparagraph free-standing
\makeatletter
\ifx\paragraph\undefined\else
  \let\oldparagraph\paragraph
  \renewcommand{\paragraph}{
    \@ifstar
      \xxxParagraphStar
      \xxxParagraphNoStar
  }
  \newcommand{\xxxParagraphStar}[1]{\oldparagraph*{#1}\mbox{}}
  \newcommand{\xxxParagraphNoStar}[1]{\oldparagraph{#1}\mbox{}}
\fi
\ifx\subparagraph\undefined\else
  \let\oldsubparagraph\subparagraph
  \renewcommand{\subparagraph}{
    \@ifstar
      \xxxSubParagraphStar
      \xxxSubParagraphNoStar
  }
  \newcommand{\xxxSubParagraphStar}[1]{\oldsubparagraph*{#1}\mbox{}}
  \newcommand{\xxxSubParagraphNoStar}[1]{\oldsubparagraph{#1}\mbox{}}
\fi
\makeatother


\usepackage{longtable,booktabs,array}
\newcounter{none} % for unnumbered tables
\usepackage{calc} % for calculating minipage widths
% Correct order of tables after \paragraph or \subparagraph
\usepackage{etoolbox}
\makeatletter
\patchcmd\longtable{\par}{\if@noskipsec\mbox{}\fi\par}{}{}
\makeatother
% Allow footnotes in longtable head/foot
\IfFileExists{footnotehyper.sty}{\usepackage{footnotehyper}}{\usepackage{footnote}}
\makesavenoteenv{longtable}
\usepackage{graphicx}
\makeatletter
\newsavebox\pandoc@box
\newcommand*\pandocbounded[1]{% scales image to fit in text height/width
  \sbox\pandoc@box{#1}%
  \Gscale@div\@tempa{\textheight}{\dimexpr\ht\pandoc@box+\dp\pandoc@box\relax}%
  \Gscale@div\@tempb{\linewidth}{\wd\pandoc@box}%
  \ifdim\@tempb\p@<\@tempa\p@\let\@tempa\@tempb\fi% select the smaller of both
  \ifdim\@tempa\p@<\p@\scalebox{\@tempa}{\usebox\pandoc@box}%
  \else\usebox{\pandoc@box}%
  \fi%
}
% Set default figure placement to htbp
\def\fps@figure{htbp}
\makeatother



\ifLuaTeX
\usepackage[bidi=basic,shorthands=off]{babel}
\else
\usepackage[bidi=default,shorthands=off]{babel}
\fi
\ifPDFTeX
\else
\babelfont{rm}[]{Fira Sans}
\fi
\ifLuaTeX
  \usepackage{selnolig} % disable illegal ligatures
\fi


\setlength{\emergencystretch}{3em} % prevent overfull lines

\providecommand{\tightlist}{%
  \setlength{\itemsep}{0pt}\setlength{\parskip}{0pt}}



 


\usepackage{tikz}
\usepackage{graphicx}
\usepackage{amsmath}
\renewcommand{\arraystretch}{1.4}
\makeatletter
\@ifpackageloaded{caption}{}{\usepackage{caption}}
\AtBeginDocument{%
\ifdefined\contentsname
  \renewcommand*\contentsname{Tabla de contenidos}
\else
  \newcommand\contentsname{Tabla de contenidos}
\fi
\ifdefined\listfigurename
  \renewcommand*\listfigurename{Listado de Figuras}
\else
  \newcommand\listfigurename{Listado de Figuras}
\fi
\ifdefined\listtablename
  \renewcommand*\listtablename{Listado de Tablas}
\else
  \newcommand\listtablename{Listado de Tablas}
\fi
\ifdefined\figurename
  \renewcommand*\figurename{Figura}
\else
  \newcommand\figurename{Figura}
\fi
\ifdefined\tablename
  \renewcommand*\tablename{Tabla}
\else
  \newcommand\tablename{Tabla}
\fi
}
\@ifpackageloaded{float}{}{\usepackage{float}}
\floatstyle{ruled}
\@ifundefined{c@chapter}{\newfloat{codelisting}{h}{lop}}{\newfloat{codelisting}{h}{lop}[chapter]}
\floatname{codelisting}{Listado}
\newcommand*\listoflistings{\listof{codelisting}{Listado de Listados}}
\makeatother
\makeatletter
\makeatother
\makeatletter
\@ifpackageloaded{caption}{}{\usepackage{caption}}
\@ifpackageloaded{subcaption}{}{\usepackage{subcaption}}
\makeatother
\usepackage{bookmark}
\IfFileExists{xurl.sty}{\usepackage{xurl}}{} % add URL line breaks if available
\urlstyle{same}
% fallback for those not using the hyperref driver hyperxmp:
\makeatletter
\@ifundefined{xmpquote}{\newcommand{\xmpquote}[1]{#1}}{}
\makeatother
\hypersetup{
  pdftitle={Práctica 3 -- Transporte a través del Xilema},
  pdflang={es},
  colorlinks=true,
  linkcolor={blue},
  filecolor={Maroon},
  citecolor={Blue},
  urlcolor={Blue},
  pdfcreator={LaTeX via pandoc}}


\title{\vspace{-1.5em}{\Huge \textbf{Práctica 3 -- Transporte a través
del Xilema}}\vspace{-1.5em}}
\date{}
\begin{document}
\maketitle


\setstretch{1.2}
\textbf{Fecha}: \hrulefill \hfill \textbf{Grupo}: \hrulefill

\textbf{Nombre:} \hrulefill

\textbf{Objetivos:}

\begin{itemize}
\tightlist
\item
  Demostrar el transporte de agua y solutos a través de los elementos
  conductores del xilema.
\item
  Evaluar el efecto de la tasa de transpiración y la cavitación/embolia
  en la distancia total de ascenso xilemático.
\end{itemize}

\subsection{1. Materiales}\label{materiales}

\begin{itemize}
\tightlist
\item
  Tallos o pecíolos frescos con hojas de apio (\emph{Apium graveolens}).
\item
  Kit de disección.
\item
  Regla graduada.
\item
  Bolsas plásticas transparentes.
\end{itemize}

\subsection{2. Aspectos Teóricos}\label{aspectos-teuxf3ricos}

El transporte de agua a larga distancia en las plantas ocurre a través
del \textbf{xilema} impulsado principalmente por la \textbf{Teoría de la
Cohesión-Tensión}. La pérdida de agua por transpiración en los estomas
genera un potencial hídrico negativo (\(\Psi_w\)) en el mesófilo,
creando una fuerza de tensión que ``jala'' la columna continua de agua
desde las raíces hasta las hojas gracias a las propiedades de cohesión y
adhesión de las moléculas de agua.

Defina brevemente:

\subsubsection{\texorpdfstring{\textbf{1. Teoría de la Cohesión-Tensión
y gradiente de
\(\Psi_w\):}}{1. Teoría de la Cohesión-Tensión y gradiente de \textbackslash Psi\_w:}}\label{teoruxeda-de-la-cohesiuxf3n-tensiuxf3n-y-gradiente-de-psi_w}

\hrulefill

\hrulefill

\hrulefill

\hrulefill

\hrulefill

\hrulefill

\subsubsection{\texorpdfstring{\textbf{2. Diferencias entre Traqueidas y
Elementos del
Vaso:}}{2. Diferencias entre Traqueidas y Elementos del Vaso:}}\label{diferencias-entre-traqueidas-y-elementos-del-vaso}

\hrulefill

\hrulefill

\hrulefill

\hrulefill

\hrulefill

\hrulefill

\subsubsection{\texorpdfstring{\textbf{3. Cavitación y Embolia
Xilemática (Causas y
efectos):}}{3. Cavitación y Embolia Xilemática (Causas y efectos):}}\label{cavitaciuxf3n-y-embolia-xilemuxe1tica-causas-y-efectos}

\hrulefill

\hrulefill

\hrulefill

\hrulefill

\hrulefill

\hrulefill

\subsection{3. Procedimiento
Experimental}\label{procedimiento-experimental}

\subsubsection{A. Tasa de Ascenso y Transpiración (Tratamientos 1 y
2)}\label{a.-tasa-de-ascenso-y-transpiraciuxf3n-tratamientos-1-y-2}

\begin{enumerate}
\def\labelenumi{\arabic{enumi}.}
\tightlist
\item
  Seleccione 6 tallos/pecíolos de tamaño homogéneo con follaje.
\item
  Realice un corte biselado bajo el agua para evitar cavitación.
\item
  Coloque los tallos en tubos con solución de colorante y deje actuar el
  sistema durante \textbf{60 minutos} ininterrumpidos bajo dos
  condiciones ambientales:

  \begin{itemize}
  \tightlist
  \item
    \textbf{Tratamiento 1 (Alta Transpiración - Luz / Aire):} 3 tallos
    (Réplicas A, B, C) expuestos a luz directa y flujo de aire.
  \item
    \textbf{Tratamiento 2 (Baja Transpiración - Oscuridad / Humedad):} 3
    tallos (Réplicas A, B, C) cubiertos con bolsa plástica en la
    oscuridad.
  \end{itemize}
\end{enumerate}

\subsubsection{B. Efecto del Bloqueo Vascular y Cavitación (Tratamiento
3)}\label{b.-efecto-del-bloqueo-vascular-y-cavitaciuxf3n-tratamiento-3}

\begin{enumerate}
\def\labelenumi{\arabic{enumi}.}
\tightlist
\item
  En 3 tallos adicionales, induzca cavitación cortándolos al aire y
  dejándolos expuestos 10 min antes de sumergir en colorante. Mantenga
  el tratamiento durante \textbf{60 minutos}.
\end{enumerate}

\begin{center}
\begin{tikzpicture}[scale=0.9]
  % ==========================================
  % TRATAMIENTO 1: Alta Transpiración
  % ==========================================
  \begin{scope}[shift={(0,0)}]
    % Recipiente y colorante
    \draw[thick] (0, 1.3) -- (0, 0.1) arc (180:360:0.5) -- (1.0, 1.3);
    \fill[blue!50!black, opacity=0.6] (0.02, 0.1) arc (180:360:0.48) -- (0.98, 0.55) -- (0.02, 0.55) -- cycle;
    \draw[blue!70!black, thick] (0.02, 0.55) -- (0.98, 0.55);

    % 3 Tallos de apio (A, B, C)
    \draw[line width=1.5pt, green!40!black] (0.25, 0.2) .. controls (0.2, 1.2) .. (0.15, 2.0);
    \draw[line width=1.5pt, green!40!black] (0.50, 0.2) -- (0.50, 2.2);
    \draw[line width=1.5pt, green!40!black] (0.75, 0.2) .. controls (0.8, 1.2) .. (0.85, 2.0);

    % Hojas / Follaje
    \fill[green!50!black] (0.15, 2.0) circle (3.5pt); \fill[green!60!black] (0.05, 2.1) circle (2.5pt); \fill[green!55!black] (0.22, 2.15) circle (2.5pt);
    \fill[green!50!black] (0.50, 2.2) circle (4.0pt); \fill[green!60!black] (0.38, 2.3) circle (3.0pt); \fill[green!55!black] (0.62, 2.3) circle (3.0pt);
    \fill[green!50!black] (0.85, 2.0) circle (3.5pt); \fill[green!60!black] (0.78, 2.15) circle (2.5pt); \fill[green!55!black] (0.95, 2.1) circle (2.5pt);

    % Etiquetas
    \node[font=\fontsize{6}{7}\bfseries\selectfont, align=center] at (0.5, -0.6) {T1: Alta Transp.};
    \node[font=\fontsize{5}{6}\selectfont, gray!50!black, align=center] at (0.5, -0.9) {Luz / Aire ($n=3$)};
  \end{scope}

  % ==========================================
  % TRATAMIENTO 2: Baja Transpiración
  % ==========================================
  \begin{scope}[shift={(2.6,0)}]
    % Recipiente y colorante
    \draw[thick] (0, 1.3) -- (0, 0.1) arc (180:360:0.5) -- (1.0, 1.3);
    \fill[blue!50!black, opacity=0.6] (0.02, 0.1) arc (180:360:0.48) -- (0.98, 0.55) -- (0.02, 0.55) -- cycle;
    \draw[blue!70!black, thick] (0.02, 0.55) -- (0.98, 0.55);

    % 3 Tallos de apio (A, B, C)
    \draw[line width=1.5pt, green!40!black] (0.25, 0.2) .. controls (0.2, 1.2) .. (0.15, 2.0);
    \draw[line width=1.5pt, green!40!black] (0.50, 0.2) -- (0.50, 2.2);
    \draw[line width=1.5pt, green!40!black] (0.75, 0.2) .. controls (0.8, 1.2) .. (0.85, 2.0);

    % Hojas / Follaje
    \fill[green!50!black] (0.15, 2.0) circle (3.5pt); \fill[green!60!black] (0.05, 2.1) circle (2.5pt); \fill[green!55!black] (0.22, 2.15) circle (2.5pt);
    \fill[green!50!black] (0.50, 2.2) circle (4.0pt); \fill[green!60!black] (0.38, 2.3) circle (3.0pt); \fill[green!55!black] (0.62, 2.3) circle (3.0pt);
    \fill[green!50!black] (0.85, 2.0) circle (3.5pt); \fill[green!60!black] (0.78, 2.15) circle (2.5pt); \fill[green!55!black] (0.95, 2.1) circle (2.5pt);

    % Bolsa plástica envolvente
    \draw[dashed, blue!70!black, fill=blue!10, fill opacity=0.35, rounded corners=4pt] 
      (-0.15, 0.8) rectangle (1.15, 2.5);
    \node[font=\fontsize{4.5}{5.5}\bfseries\selectfont, blue!80!black] at (0.5, 2.65) {Bolsa plástica};

    % Etiquetas
    \node[font=\fontsize{6}{7}\bfseries\selectfont, align=center] at (0.5, -0.6) {T2: Baja Transp.};
    \node[font=\fontsize{5}{6}\selectfont, gray!50!black, align=center] at (0.5, -0.9) {Bolsa ($n=3$)};
  \end{scope}

  % ==========================================
  % TRATAMIENTO 3: Cavitación / Bloqueo
  % ==========================================
  \begin{scope}[shift={(5.2,0)}]
    % Recipiente y colorante
    \draw[thick] (0, 1.3) -- (0, 0.1) arc (180:360:0.5) -- (1.0, 1.3);
    \fill[blue!50!black, opacity=0.6] (0.02, 0.1) arc (180:360:0.48) -- (0.98, 0.55) -- (0.02, 0.55) -- cycle;
    \draw[blue!70!black, thick] (0.02, 0.55) -- (0.98, 0.55);

    % 3 Tallos de apio (A, B, C)
    \draw[line width=1.5pt, green!40!black] (0.25, 0.2) .. controls (0.2, 1.2) .. (0.15, 2.0);
    \draw[line width=1.5pt, green!40!black] (0.50, 0.2) -- (0.50, 2.2);
    \draw[line width=1.5pt, green!40!black] (0.75, 0.2) .. controls (0.8, 1.2) .. (0.85, 2.0);

    % Bloqueo / Sello de vaselina en la base de los tallos
    \fill[orange!90!black] (0.18, 0.25) rectangle (0.32, 0.40);
    \fill[orange!90!black] (0.43, 0.25) rectangle (0.57, 0.40);
    \fill[orange!90!black] (0.68, 0.25) rectangle (0.82, 0.40);

    % Hojas / Follaje
    \fill[green!50!black] (0.15, 2.0) circle (3.5pt); \fill[green!60!black] (0.05, 2.1) circle (2.5pt); \fill[green!55!black] (0.22, 2.15) circle (2.5pt);
    \fill[green!50!black] (0.50, 2.2) circle (4.0pt); \fill[green!60!black] (0.38, 2.3) circle (3.0pt); \fill[green!55!black] (0.62, 2.3) circle (3.0pt);
    \fill[green!50!black] (0.85, 2.0) circle (3.5pt); \fill[green!60!black] (0.78, 2.15) circle (2.5pt); \fill[green!55!black] (0.95, 2.1) circle (2.5pt);

    % Etiquetas
    \node[font=\fontsize{6}{7}\bfseries\selectfont, align=center] at (0.5, -0.6) {T3: Cavitación};
    \node[font=\fontsize{5}{6}\selectfont, gray!50!black, align=center] at (0.5, -0.9) {Luz / Aire ($n=3$)};
  \end{scope}
\end{tikzpicture}
\end{center}

\subsubsection{C. Medición y Cuantificación del Ascenso del
Colorante}\label{c.-mediciuxf3n-y-cuantificaciuxf3n-del-ascenso-del-colorante}

\begin{enumerate}
\def\labelenumi{\arabic{enumi}.}
\tightlist
\item
  \textbf{Medición del ascenso del colorante:}
\end{enumerate}

\begin{itemize}
\tightlist
\item
  Coloque el pecíolo sobre una superficie plana y limpia.
\item
  Realice cortes transversales en el peciolo del apio cada 1 cm desde la
  base.
\item
  Siga realizando cortes hasta que los haces del xilema no se observen
  teñidos.
\item
  Determine la distancia recorrida por el colorante en los peciolos de
  apio de cada tratamiento.
\item
  Calcule la velocidad promedio de ascenso al minuto 60
  (\(v = \text{distancia}/60\text{ min}\)) y determine promedio
  (\(\bar{x}\)) y desviación estándar (\(s\)) para cada grupo:
\end{itemize}

\[\bar{x} = \frac{\sum_{i=1}^{n} x_i}{n} \qquad s = \sqrt{\frac{\sum_{i=1}^{n} (x_i - \bar{x})^2}{n - 1}}\]

\begin{enumerate}
\def\labelenumi{\arabic{enumi}.}
\setcounter{enumi}{1}
\tightlist
\item
  \textbf{Análisis Anatómico (Corte Transversal):}
\end{enumerate}

\begin{itemize}
\tightlist
\item
  Realice cortes transversales muy delgados a 2 cm, 5 cm y 10 cm de la
  base.
\item
  Monte en portaobjetos, observe al microscopio y determine si la
  tinción es exclusiva del xilema o si hubo difusión al parénquima
  adyacente.
\end{itemize}

\subsection{4. Registro de Resultados}\label{registro-de-resultados}

\subsubsection{Tabla 1. Altura de Ascenso y
Velocidad}\label{tabla-1.-altura-de-ascenso-y-velocidad}

\begin{center}
\resizebox{\linewidth}{!}{%
\begin{tabular}{l|l|p{2em}|p{2em}|p{2em}|p{2em}|p{2em}}
\hline
 &  & \multicolumn{3}{c|}{\textbf{Réplicas}} & \textbf{Prom} & \textbf{DE} \\
\textbf{Tratamiento} & \textbf{Variable} & \textbf{A} & \textbf{B} & \textbf{C} & \textbf{$\bar{x}$} & \textbf{$s$} \\
\hline
1. Alta Transp. & Altura & & & & & \\[1ex]
\cline{2-7}
 & Velocidad & & & & & \\[1ex]
\hline
2. Baja Transp. & Altura & & & & & \\[1ex]
\cline{2-7}
 & Velocidad & & & & & \\[1ex]
\hline
3. Cavitación  & Altura & & & & & \\[1ex]
\cline{2-7}
 & Velocidad & & & & & \\[1ex]
\hline
\end{tabular}%
}
\end{center}

Interprete los resultados de la altura de ascenso y velocidad:

\hrulefill

\hrulefill

\hrulefill

\hrulefill

\subsubsection{Tabla 2. Observación Anatómica del Corte
Transversal}\label{tabla-2.-observaciuxf3n-anatuxf3mica-del-corte-transversal}

\begin{center}
\resizebox{\linewidth}{!}{%
\begin{tabular}{l|p{14em}}
\hline
\textbf{Tratamiento} & \textbf{Observación Microscópica / Disposición de Haces Teñidos} \\
\hline
1. Alta Transp. &  \\[20pt]
\hline
2. Baja Transp. & \\[20pt]
\hline
3. Cavitación & \\[20pt]
\hline
\end{tabular}%
}
\end{center}

Interprete los resultados de la observación anatómica:

\hrulefill

\hrulefill

\hrulefill

\hrulefill

\subsection{5. Análisis Gráfico}\label{anuxe1lisis-gruxe1fico}

Represente la altura final y velocidad de ascenso utilizando gráficas
con puntos (promedio) y barras de error (desviación estándar) por
tratamiento.

\subsubsection{Gráfica 1: Altura Final de
Ascenso}\label{gruxe1fica-1-altura-final-de-ascenso}

\begin{center}
\begin{tikzpicture}[scale=0.95]
  % Grid ggplot2
  \draw[gray!20, step=0.5cm] (0,0) grid (6,4.0);
  \draw[gray!50, step=1cm] (0,0) grid (6,4.0);
  \draw[thick] (0,0) rectangle (6,4.0);
  
  % Eje X
  \node[below, font=\tiny] at (1,0) {\textbf{Alta Transp.}};
  \node[below, font=\tiny] at (3,0) {\textbf{Baja Transp.}};
  \node[below, font=\tiny] at (5,0) {\textbf{Cavitación}};
  \node[below=10pt, font=\small] at (3.25,0) {\textbf{Tratamiento Experimental}};
  
  % Eje Y
  \node[rotate=90, above=10pt, font=\small] at (0,2.0) {\textbf{Altura de Ascenso (cm)}};
\end{tikzpicture}
\end{center}

Interprete los resultados observados:

\hrulefill

\hrulefill

\hrulefill

\hrulefill

\hrulefill

\subsubsection{Gráfica 2: Velocidad de Ascenso
Xilemático}\label{gruxe1fica-2-velocidad-de-ascenso-xilemuxe1tico}

\begin{center}
\begin{tikzpicture}[scale=0.95]
  % Grid ggplot2
  \draw[gray!20, step=0.5cm] (0,0) grid (6,4.0);
  \draw[gray!50, step=1cm] (0,0) grid (6,4.0);
  \draw[thick] (0,0) rectangle (6,4.0);
  
  % Eje X
  \node[below, font=\tiny] at (1,0) {\textbf{Alta Transp.}};
  \node[below, font=\tiny] at (3,0) {\textbf{Baja Transp.}};
  \node[below, font=\tiny] at (5,0) {\textbf{Cavitación}};
  \node[below=10pt, font=\small] at (3.25,0) {\textbf{Tratamiento Experimental}};
  
  % Eje Y
  \node[rotate=90, above=10pt, font=\small] at (0,2.0) {\textbf{Velocidad ($\text{cm/min}$)}};
\end{tikzpicture}
\end{center}

Interprete los resultados observados:

\hrulefill

\hrulefill

\hrulefill

\hrulefill

\hrulefill

\subsection{6. Actividades
Complementarias}\label{actividades-complementarias}

\begin{enumerate}
\def\labelenumi{\arabic{enumi}.}
\tightlist
\item
  Explique el mecanismo anatómico y biofísico por el cual los cortes
  realizados al aire sufren una reducción significativa en la tasa de
  transporte xilemático.
\end{enumerate}

\hrulefill

\hrulefill

\hrulefill

\hrulefill

\hrulefill

\hrulefill

\hrulefill

\hrulefill

\begin{enumerate}
\def\labelenumi{\arabic{enumi}.}
\setcounter{enumi}{1}
\tightlist
\item
  ¿Por qué la coloración se restringe a los haces vasculares en el corte
  transversal y no se difunde uniformemente por todo el tejido
  parenquimático al finalizar el experimento?
\end{enumerate}

\hrulefill

\hrulefill

\hrulefill

\hrulefill

\hrulefill

\hrulefill

\hrulefill

\hrulefill

\begin{enumerate}
\def\labelenumi{\arabic{enumi}.}
\setcounter{enumi}{2}
\tightlist
\item
  ¿Qué mecanismos fisiológicos emplean las plantas para reparar o
  recuperar la funcionalidad de los vasos cavitados?
\end{enumerate}

\hrulefill

\hrulefill

\hrulefill

\hrulefill

\hrulefill

\hrulefill

\hrulefill

\hrulefill

\subsection{7. Conclusión}\label{conclusiuxf3n}

Relacione los resultados cuantitativos con la tensión transpiratoria e
interprete la importancia del xilema como vía principal de transporte
hídrico:

\hrulefill

\hrulefill

\hrulefill

\hrulefill

\hrulefill

\hrulefill

\hrulefill

\subsection{8. Bibliografía}\label{bibliografuxeda}

\begin{itemize}
\tightlist
\item
  Azcón-Bieto, J., \& Talón, M. (2008). \emph{Fundamentos de Fisiología
  Vegetal} (2.ª ed.). McGraw-Hill.
\item
  Melgarejo, L. M. (Ed.). (2010). \emph{Experimentos en fisiología
  vegetal}. Universidad Nacional de Colombia.
\item
  Taiz, L., \& Zeiger, E. (2002). \emph{Plant Physiology} (3rd ed.).
  Sinauer Associates.
\end{itemize}




\end{document}
